
\section{Robustness and Extensions}
\label{sec:robustness}


We subject the baseline to five groups of checks:
(i)~FL-definition variants (IQR multiplier, win-rate cutoff, temporal
window); (ii)~sample and specification variants (unrestricted sample,
homogeneous subsamples, item$\times$year FE, matching);
(iii)~alternative clustering; (iv)~sensitivity to unobservables
(Cinelli--Hazlett); and (v)~staggered DiD.
Table~\ref{tab:robustness_summary} collects the point estimates, with
an extra column spelling out what each estimator assumes.


\begin{table}[htbp]
\centering
\caption{Robustness summary: FL conditional price coefficient across checks}
\label{tab:robustness_summary}
\begin{adjustbox}{max width=\textwidth}
\footnotesize
\begin{tabular}{lcccp{3.2cm}l}
\hline\hline
Check & Coeff. & SE & $N$ & Assumption & Note \\
\hline
\multicolumn{6}{l}{\textit{Baseline and matching}} \\
OLS (item + year + PBU FE) & 0.064 & 0.020 & 1{,}654{,}401 & Sel.\ on obs.\ + FE & Baseline \\
CEM & 0.077 & 0.024 & 969{,}751 & Sel.\ on obs. & Sec.~\ref{sec:results_ols} \\
IPW & 0.055 & 0.021 & 830{,}194 & Sel.\ on obs. & Sec.~\ref{sec:results_ols} \\
Cross-fit & 0.036 & 0.019 & 1{,}654{,}401 & Sel.\ on obs.\ + FE; no mechanical link & Out-of-sample \\
\hline
\multicolumn{6}{l}{\textit{FL definition variants}} \\
IQR 1.0$\times$ & 0.079 & 0.023 & 1{,}654{,}401 & Sel.\ on obs.\ + FE & 3{,}442 FL \\
IQR 1.5$\times$ & \multicolumn{5}{l}{\textit{(= OLS baseline above)}} \\
IQR 2.0$\times$ & 0.060 & 0.022 & 1{,}654{,}401 & Sel.\ on obs.\ + FE & 2{,}093 FL \\
IQR 3.0$\times$ & 0.050 & 0.024 & 1{,}654{,}401 & Sel.\ on obs.\ + FE & 1{,}456 FL \\
\hline
\multicolumn{6}{l}{\textit{Specification variants}} \\
Item$\times$year FE & 0.074 & 0.022 & 1{,}654{,}401 & Sel.\ on obs.\ + FE & Tighter controls \\
Two-way clustering & 0.064 & 0.024 & 1{,}654{,}401 & Sel.\ on obs.\ + FE & Item + PBU \\
\hline
\multicolumn{6}{l}{\textit{Alternative specifications}} \\
Continuous (log max AL) & 0.022 & 0.005 & 1{,}654{,}401 & Sel.\ on obs.\ + FE & Per log-point \\
Decomposition (odd-yr FL) & 0.043 & 0.019 & 1{,}654{,}401 & Sel.\ on obs.\ + FE & Intermediate \\
Horse-race (FL + Imhof CV) & 0.084 & 0.023 & 1{,}654{,}401 & Sel.\ on obs.\ + FE & Suppression \\
IV (leave-one-out) & 0.194 & 0.077 & 1{,}654{,}401 & Excl.\ restriction & Meas.\ error \\
\hline
\multicolumn{6}{l}{\textit{Staggered DiD}} \\
Stacked DiD & $-$0.006 & 0.014 & 715{,}116 & Parallel trends & Pre-trends \\
Oster $\hat{\delta}$ & \multicolumn{2}{c}{degen.} & 1{,}654{,}401 & Proportional sel. & $R^2$ gap $\approx 0$ \\
\hline\hline
\end{tabular}
\end{adjustbox}
\begin{tablenotes}
\footnotesize
\item \textit{Notes:} ``Sel.\ on obs.\ + FE'' = selection on observables absorbed by item, year, and PBU fixed effects. Cross-fit removes any mechanical link between classification and outcome. The IV is a measurement-error diagnostic; exclusion-restriction concerns (Section~\ref{sec:results_ols}) keep it off the primary range.
\end{tablenotes}
\end{table}


\paragraph{Threshold sensitivity.}
The coefficient decreases monotonically with stricter thresholds
(0.079 at 1.0$\times$, 0.064 baseline, 0.060 at 2.0$\times$, 0.050
at 3.0$\times$), most likely reflecting classification noise as the
treated group shrinks. A two-dimensional sensitivity analysis varying
both multiplier and win-rate cutoff produces significant coefficients
across all 36 cells
(Figure~\ref{fig:threshold_heatmap}, Appendix~\ref{app:robustness}).
For a deploying authority, this means the screen's triage signal is
not sensitive to the exact threshold chosen.

\paragraph{Sensitivity to unobservables.}
\citet{cinelli2020making} sensitivity analysis yields $RV_{q=1} =
17.5\%$: an unobserved confounder would need to explain 17.5\% of
residual variation in \emph{both} FL presence and log prices to nullify
the coefficient---stringent given the item, year, and PBU fixed-effects
structure. The bound assumes a single additive confounder; if selection
operates through multiple correlated channels---each individually below
17.5\%---the required confounder strength may be lower. The positive
pre-trends at $t = -2$ and $t = -3$ suggest multi-channel selection
may be relevant (Section~\ref{sec:limitations}). Full sensitivity
plots are in Appendix~\ref{app:robustness}.


\paragraph{Staggered difference-in-differences.}
The Callaway \& Sant'Anna ATT is 0.014 (SE~$= 0.042$) and the stacked
DiD \citep{cengiz2019effect} gives $-0.006$ (SE~$= 0.014$): both are
statistically indistinguishable from zero. The minimum detectable effect
at 80\% power is 0.117 for the C\&S estimator---nearly twice the OLS
baseline of 0.064---so the DiD lacks statistical power to detect the
cross-sectional association even if it were entirely causal. Moreover,
positive pre-event coefficients at $t = -2$ and $t = -3$
(Figure~\ref{fig:event_study}, Appendix~\ref{app:did}) indicate that
PBUs receiving FL entry already had higher prices, consistent with
the framework's prediction that cover bidders target profitable
environments.

The DiD null is uninformative about the FL--price association: it is
compatible with zero causal effect, insufficient power, or---most
likely---violation of parallel trends through strategic market
selection. The paper's evidence hierarchy does not rely on DiD
identification; the cross-sectional association documented in
Section~\ref{sec:results_ols} rests on selection-on-observables
absorbed by high-dimensional fixed effects
(Tables~\ref{tab:stacked_did}--\ref{tab:did_revised},
Appendix~\ref{app:did}).

\paragraph{Coefficient stability (Oster).}
The \citet{oster2019unobservable} bound is degenerate
($\hat{\delta} = 261.6$) because PBU FE barely move $R^2$
(0.879 to 0.886)---a design strength, not fragility. The
\citet{cinelli2020making} framework ($RV = 17.5\%$) is the
appropriate sensitivity metric
(Table~\ref{tab:oster_delta}, Appendix~\ref{app:robustness}).


\paragraph{Cost-benefit of screening.}
Appendix~\ref{app:extensions} reports cost-benefit calculations for
screening deployment, applying the price coefficients to FL-present
spending under varying assumptions about investigation costs and
deterrence rates.


\paragraph{Oversight heterogeneity.}
The FL--price association varies sharply with PBU size: 0.214 in the
smallest quartile, 0.098 in Q2, 0.045 in Q3, and 0.017 in Q4---a
12.5$\times$ gradient
(Table~\ref{tab:regime_oversight}, Appendix~\ref{app:robustness}).
This matches the framework's comparative static
($\partial m^*/\partial \theta_k < 0$): where oversight is weaker,
the screen bites harder. Smaller PBUs also have thinner fixed-effect
cells, so part of the gradient could be mechanical, but the monotonic
decline across all four quartiles is hard to attribute to noise
alone.
Threshold stability across IQR multipliers and a comparison between
the IQR rule and model-based classifiers are in
Appendix~\ref{app:robustness}.
